% Appendix A: full proofs (r470). Included from paper.tex.
\section{Proofs}
\label{app:proofs}

\subsection{Monotonicity of the fixed-$k$ flip probability (Lemma~\ref{lem:exact})}

Recall
\[
f_K(k)=\P\big(\MV(k)\neq\MV(N)\,\big|\,K\big),
\qquad x_k\sim\HG(N,K,k),
\]
with the negative side winning ties ($\MV(k)=1$ iff $x_k>k/2$).
We prove that $f_K(k)$ is nonincreasing along the odd numbers
$k\in\{1,3,5,\dots,N-1\}$ for every fixed $K$; this suffices for the paper,
since its early-budget candidates are odd. For a selected odd subgrid,
$k_H^{\mathrm{odd}}(\alpha)$ is defined only when its feasible set
$\mathcal F_H(\alpha)=\{k:\E_H[f_K(k)]\le\alpha\}$ is nonempty; otherwise
the protocol uses its separate full-$N$ fallback rather than calling it an odd
certificate.

\paragraph{Setup and case split.}
$\MV(k)\neq\MV(N)$ can happen in only two ways:
\begin{itemize}
\item \textbf{Case U (upset).} $K<N/2$ (full vote negative) and $x_k>k/2$
(prefix votes positive).
\item \textbf{Case R (reversal).} $K>N/2$ (full vote positive) and $x_k\le k/2$
(prefix votes negative).
\end{itemize}
When $N$ is even there is also the center case $K=N/2$, handled separately
below. Write $f_K(k)=U_K(k)+R_K(k)$ with $U_K(k)=\P(x_k>k/2)\,\mathbf 1\{K<N/2\}$
and $R_K(k)=\P(x_k\le k/2)\,\mathbf 1\{K>N/2\}$. At most one of the two terms
is nonzero away from the center.

\paragraph{The reflection identity.}
Swapping the labels ``pass''/``fail'' of all $N$ rollouts maps the law
$\HG(N,K,k)$ of $x_k$ onto the law of $k-x_k$ under count $N-K$, i.e.\
$\P_{K}(x_k>s)=\P_{N-K}(x_k<k-s)$. Hence, for odd $k$,
\[
U_K(k)=\P_K(x_k>k/2)=\P_{N-K}(x_k<k/2)=R_{N-K}(k)
\qquad(K<N/2),
\]
since $k$ odd makes $k/2$ a half-integer, so the strict and weak
inequalities coincide. Every upset probability thus equals a reversal
probability at the reflected count, and it suffices to prove
\begin{equation}
\label{eq:revmain}
R_K(k+2)\le R_K(k)\qquad\text{for all }K>N/2\text{ and odd }k,
\end{equation}
which covers Case R directly and Case U by reflection (the same display
applied at count $N-K$).

\paragraph{Even-$N$ center case.}
If $N$ is even and $K=N/2$, symmetry under $x\mapsto k-x$ and the absence of
an odd-$k$ tie give $\P(x_k>k/2)=\P(x_k<k/2)=1/2$. The full vote is negative,
so $f_{N/2}(k)=1/2$ for every odd $k$: it is constant. This is why reflection
is restricted to $K<N/2$.

\paragraph{Coupling construction.}
Fix $K>N/2$ and odd $k\le N-3$. We couple the two hypergeometric draws on a
single probability space. Let the population of $N$ rollouts consist of $K$
passes and $N-K$ fails. Draw an \emph{ordered} sample of size $k+2$ uniformly
without replacement and split it as
\[
\underbrace{X_1,\dots,X_k}_{\text{prefix }k}\;\Big|\;
\underbrace{X_{k+1},X_{k+2}}_{\text{extra pair }(A,B)} .
\]
The marginal law of $(X_1,\dots,X_k)$ is exactly $\HG(N,K,k)$, and
$x_{k+2}=x_k+A+B$ has law $\HG(N,K,k+2)$. All statements below are about this
joint space, so they hold for the two marginals simultaneously.

Let $t=k/2+1/2$ be the (integer) flip threshold: the prefix at $k$ votes
negative iff $x_k\le t-1$; the prefix at $k+2$ votes negative iff
$x_{k+2}\le t$. Partition the event $\{x_k=s\}$ by the value of the pair
$A+B\in\{0,1,2\}$:
\[
\begin{aligned}
E_0&=\{x_k=t-1,\ A+B=0\}, &&\text{flip at }k,\ \text{flip at }k+2,\\
E_1&=\{x_k=t-1,\ A+B=1\}, &&\text{flip at }k,\ \text{flip at }k+2,\\
E_2&=\{x_k=t-1,\ A+B=2\}, &&\text{flip at }k,\ \text{no flip at }k+2,\\
F_0&=\{x_k=t,\ \ \ A+B=0\}, &&\text{no flip at }k,\ \text{flip at }k+2,\\
F_1&=\{x_k=t,\ \ \ A+B=1\}, &&\text{no flip at }k,\ \text{no flip at }k+2.
\end{aligned}
\]
All other states agree on both sides (e.g.\ $x_k\le t-2$ flips under both, since
$A+B\le2$; $x_k\ge t+1$ flips under neither). Therefore
\begin{equation}
\label{eq:rdiff}
R_K(k)-R_K(k+2)=\P(E_2)-\P(F_0),
\end{equation}
and it remains to show $\P(E_2)\ge\P(F_0)$.

\paragraph{Comparison of the two boundary atoms, including support boundaries.}
Both atoms factor as (prefix probability) $\times$ (pair probability).  We
first split off the case $\P(F_0)=0$.  Then \eqref{eq:rdiff} gives
$R_K(k)-R_K(k+2)=\P(E_2)\ge0$ directly; in particular, no ratio is formed
when a binomial factor in $F_0$ vanishes (including a zero remaining-fail
count or an infeasible $x_k=t$ atom).  It remains to consider
$\P(F_0)>0$.  This support condition makes every denominator below strictly
positive and makes the ratio algebra legitimate. With $k=2t-1$, first,
\[
\P(x_k=t-1)=\frac{\binom{K}{t-1}\binom{N-K}{t}}{\binom{N}{k}},
\qquad
\P(x_k=t)=\frac{\binom{K}{t}\binom{N-K}{t-1}}{\binom{N}{k}},
\]
since $k-(t-1)=t$ and $k-t=t-1$. Given $x_k=t-1$, the remainder holds
$K-t+1$ passes out of $N-k$, so
$\P(A+B=2\mid x_k=t-1)=\binom{K-t+1}{2}\big/\binom{N-k}{2}$; given $x_k=t$,
the remainder holds $N-K-t+1$ fails, so
$\P(A+B=0\mid x_k=t)=\binom{N-K-t+1}{2}\big/\binom{N-k}{2}$. The
denominators $\binom{N}{k}\binom{N-k}{2}$ are common to both atoms. Under
$\P(F_0)>0$, hence
\[
\frac{\P(E_2)}{\P(F_0)}
=\frac{\binom{K}{t-1}\binom{N-K}{t}\binom{K-t+1}{2}}
{\binom{K}{t}\binom{N-K}{t-1}\binom{N-K-t+1}{2}}
=\frac{K-t}{N-K-t},
\]
where the last equality uses the three elementary ratios
$\binom{K}{t-1}\big/\binom{K}{t}=\frac{t}{K-t+1}$,
$\binom{N-K}{t}\big/\binom{N-K}{t-1}=\frac{N-K-t+1}{t}$, and
$\binom{K-t+1}{2}\big/\binom{N-K-t+1}{2}=\frac{(K-t+1)(K-t)}{(N-K-t+1)(N-K-t)}$,
whose product telescopes to the displayed ratio. Therefore, with
$C:=\binom{K}{t}\binom{N-K}{t-1}\binom{N-K-t+1}{2}\big/
\big[\binom{N}{k}\binom{N-k}{2}\big]\ge 0$,
\begin{equation}
\label{eq:atomdiff}
\P(E_2)-\P(F_0)=C\big[(K-t)-(N-K-t)\big]=C\,(2K-N)\;\ge\;0,
\end{equation}
because $K>N/2$ in Case R.  The excluded support case was already proved
directly; if $t>K$, $\P(x_k=t-1)=0$ and both atoms vanish. By \eqref{eq:rdiff}
this proves \eqref{eq:revmain} and, by the
reflection identity, $U_K(k+2)\le U_K(k)$ as well. Together with the
even-$N$ center case, this completes the proof of monotonicity along odd
$k$.\qed

\begin{remark}[Verification at $N=32$]
Independently of the proof, the revision's exact-rational CPU audit enumerates
all admissible $(N,K,k)$ for $2\le N\le64$, including the boundary
$(N,K,k)=(32,17,29)$, checks \eqref{eq:rdiff}, the support split, and the
monotonicity inequality; every even-$N$ center is $1/2$ on the odd grid, and
a center point mass with $\alpha<1/2$ invokes the explicit full-$N$ fallback.
The historical $N=32$ artifact remains a separate frozen result check.
\end{remark}

\begin{remark}[Why only odd $k$ enters the certificate]
Requiring $k$ odd removes tie-breaking from the certificate family
(every prefix has a strict majority side). Its binary search is only over a
nonempty feasible odd ladder; FULL-$N$ is a separate zero-flip fallback.
Nothing else in the paper uses monotonicity at even $k$.
\end{remark}

\subsection{Oracle terminal-stop bookkeeping (Theorem~\ref{thm:tower})}

For the oracle theorem, the stopping rule is
$\tau=\min\{k\in\{3,\ldots,N\}:\oraclecert(x_k,k)\le\alpha\}$.  At a reachable
terminal state, all $N$ binary outcomes have been observed, hence $x_N=K$.
We define the endpoint certificate to be zero; whenever the posterior is
supported, this also follows because the hypergeometric likelihood at
$(x_N,N)$ has support only on $K=x_N$:
\[
\oraclecert(x_N,N)=0.
\]
Thus $k=N$ satisfies the same threshold rule rather than being an exceptional
terminal state. The tower argument applies to this bounded stopping time,
and
\[
\P_{\trueH}(\flip)=\E_{\trueH}\big[\oraclecert(x_\tau,\tau)\big]
=\E_{\trueH}\big[\oraclecert\,\mathbf 1\{\tau<N\}\big]+\E_{\trueH}\big[\oraclecert\,\mathbf 1\{\tau=N\}\big].
\]
The first term is at most $\alpha\,\P(\tau<N)$, and the second term is zero.
Therefore the integrand is bounded pointwise by $\alpha$ and
$\P(\flip)\le\alpha$.  For each fixed true count $K$, a path reaching $N$
also has $x_N=K$, so its terminal contribution to the DP flip probability
$g(K;\alpha,\widehat H_m)$ is exactly zero. This endpoint fact also holds
for the fitted DP, but it does not identify a fitted earlier-prefix score with
the oracle conditional probability.

\begin{remark}[The terminal-stop mass is small in practice]
On the TEST split, the exact fraction of paths reaching $k=N$ under
BAYES-H is $0.59\%$ at $\alpha=0.05$ and $1.10\%$ at $\alpha=0.02$ (the
fitted score stops early on the reported bimodal mass). Their contribution to
replay flip is zero; their contribution to the expected stopping count is
included exactly (Appendix~\ref{app:dp}).
\end{remark}

\subsection{Linearity of the population flip rate}

For the population statement, FIT and CAL are assumed independent i.i.d.
task samples from the same count-exchangeable replay population. Conditional
on FIT, every prior, fitted-score table, and candidate rule is frozen before
CAL is inspected. For such a fixed rule and fixed prior $\widehat H_{\rm FIT}$
used \emph{inside} the fitted score, the per-task flip probability $g(K)$
depends only on the true count $K$ (via the hypergeometric law of prefixes)
and on $\widehat H_{\rm FIT}$ (through the posterior threshold). The CAL
values $g(K_i)$ are consequently i.i.d. bounded observations conditional on
FIT, and the population replay flip rate is
\[
\P_{\text{replay},\trueH}(\flip)=\sum_{K=0}^{N} \trueH[K]\,g(K),
\]
a \emph{linear} functional of the true mixture $\trueH$. Conditional on FIT,
empirical Bernstein on the CAL mean bounds the population quantity for a
frozen rule, with each per-task quantity $g\in[0,1]$. Bonferroni over the
pre-declared family of $J=64$ rules permits CAL-side selection because
simultaneous UCBs cover all $J$ frozen linear functionals. The fixed benchmark
TEST split is not used in this argument and remains descriptive. This
population result is distinct from the oracle conditional identity in
Theorem~\ref{thm:tower}.
